\section{大-$\beta_0$ 系数函数}\label{sect_bubbles}%
%%%%%%%%%%%%%%%%%%%%%%%%%%%%%%%%%%%%%%%%%%%%%%%%%%%%%
%

qITD 系数函数中 renormalon 奇点的领头贡献示于图~\ref{fig:bubblechain}，
技术细节见附录 A。我们得到
\begin{align}
\label{B[H]}
B[H^\parallel](w)&=\frac{2C_F}{w} \biggl\{\Big[\frac{1+\alpha^2}{1-\alpha}
- \Big(2\alpha\,{}_2F_1(1,2-w,2+w,\alpha)
\notag\\&\quad
 +\bar \alpha (1-w^2)\Big)\alpha^w\,h_0(w,X)\Big]_+
\notag\\ &\quad
+\delta(\bar \alpha)\bigg[\frac{3(w^2-w-1)}{(w+2)(2w-1)}h_0(w,X)-\frac{3}{2}\bigg]\biggr\}
\notag\\[1mm] &\quad
+\widetilde{R}(w),
\\
B[H^\perp](w)&= B[H^\parallel](w)-4C_F \bar \alpha(1+w)\alpha^w\,h_0(w,X)\,,
\end{align}
其中 $\bar\alpha = 1-\alpha$，
\begin{align}
& h_0(w,X)=X^w \frac{\Gamma(1-w)}{\Gamma(2+w)},
\notag\\
%&G_0(w)=\frac{\Gamma(4+2w)}{6\Gamma(1-w)\Gamma(1+w)\Gamma^2(2+w)},
%\notag\\[4mm]
&X=- \frac{v^2 z^2 \mu^2 e^{5/3}}{4},
\end{align}
而函数 $\widetilde{R}(w)$ 定义为另一函数的级数展开
\begin{align}
R(w)&=2C_F \biggl\{\Big[\frac{1+\alpha^2}{1-\alpha}\frac{\alpha^wG_0(w)\!-\!1}{w}+\alpha^w \bar \alpha (2\!+\!w)G_0(w)\Big]_+
\notag\\&\quad
+\frac{\delta(\bar \alpha)}{w}\bigg[\frac{3}{2}-\frac{2w+3}{(w+2)(w+1)}G_0(w)\bigg]\biggr\},
\notag\\
G_0(w)&=\frac{\Gamma(4+2w)}{6\Gamma(1-w)\Gamma(1+w)\Gamma^2(2+w)},
\end{align}
且满足
\begin{align}
  R(w) = \sum_n w^n R_n\,, \qquad \widetilde{R}(w) = \sum_n \frac{w^n}{(n+1)!} R_n \,.
\end{align}
把 Borel 变换在 $w=0$ 处 Taylor 展开即得以耦合常数表示的系数函数微扰展开。
$\mathcal{O}(\alpha_s)$ 项为
\begin{widetext}
\begin{eqnarray}
\label{NLO}
H^{(1)\parallel}(\alpha)&=&2C_F\Bigg[
\bigg(-\mathbf{L}_\mu \frac{1+\alpha^2}{1-\alpha}+\frac{3-8\alpha+3\alpha^2-4\ln\bar \alpha}{1-\alpha}\bigg)_+
+\delta(\bar \alpha)\bigg(\frac{3}{2}\mathbf{L}_\mu+\frac{7}{2}\bigg)\Bigg],
\\
H^{(1)\perp}(\alpha)&=&2C_F\Bigg[
 \bigg(-\mathbf{L}_\mu \frac{1+\alpha^2}{1-\alpha}+\frac{1-4\alpha+\alpha^2-4\ln\bar \alpha}{1-\alpha}\bigg)_+
 +\delta(\bar \alpha)\bigg(\frac{3}{2}\mathbf{L}_\mu+\frac{5}{2}\bigg)\Bigg],
\end{eqnarray}
而 $\mathcal{O}(\alpha_s^2)$ 修正的 $n_f$ 部分为
\begin{eqnarray}
\label{NNLO}
H^{(2)\parallel}(\alpha)&=&
-\frac{4C_Fn_f}{3}\biggl\{-\frac{\mathbf{L}_\mu^2}{2}\Big(\frac{1+\alpha^2}{1-\alpha}\Big)_+-\mathbf{L}_\mu \Big[\frac{1+\alpha^2}{1-\alpha}\Big(\ln\alpha+\frac{2}{3}\Big)+4\frac{\alpha+\ln\bar \alpha}{1-\alpha}\Big]_+
+\delta(\bar \alpha)\Big(\frac{3\mathbf{L}_\mu^2}{4}+\frac{19}{4}\mathbf{L}_\mu+\frac{159}{16}\Big)
\\\nonumber && \qquad+\Big[\frac{1+\alpha^2}{1-\alpha}\Big(\frac{41}{18}+\frac{7}{6}\ln\alpha-\frac{\ln^2\alpha}{4}\Big)
-\frac{4}{1-\alpha}\Big(\text{Li}_2(\alpha)+\ln^2\bar \alpha+\frac{8}{3}\ln\bar \alpha+\ln \bar \alpha \ln\alpha\Big)
-\frac{\alpha(13+6\ln \alpha)}{1-\alpha}\Big]_+ \biggr\}.
\\
H^{(2)\perp}(\alpha)&=&H^{(1)\parallel}(\alpha)+\frac{4C_Fn_f}{3}\bar \alpha\Big(2\mathbf{L}_\mu+2\ln\alpha+\frac{10}{3}\Big),
\end{eqnarray}
\end{widetext}
其中
\begin{eqnarray}
\mathbf{L}_\mu=\ln\bigg(\frac{-v^2z^2\mu^2}{4e^{-2\gamma_E}}\bigg).
\end{eqnarray}
我们已检验：式 \eqref{NLO} 经 Fourier 变换 \eqref{def:qPDF+pPDF} 后给出文献 \cite{Xiong:2013bka,Stewart:2017tvs} 中计算的 qPDF 单圈修正。

%
\subsection{Borel 变换的奇点}\label{sect_renormalons}%
%
系数函数 Borel 变换的奇点结构如图~\ref{fig:borelplane} 所示。
%
\begin{figure}[b]
\centering
 \includegraphics[width=0.45\textwidth]{BVZplots/borelplane}
  %~\includegraphics[width=6.0cm]{BMJplots/fit-t4-qPDF}
\caption{Borel 变换的奇点结构}
\label{fig:borelplane}
\end{figure}
%
在 $w=1/2$ 处有一个紫外 renormalon 奇点，在正整数 $w=1,2,\ldots$ 处有一系列红外 renormalon。在我们计算的精度（单条气泡链）下所有奇点都是单极点。
这些奇点阻碍了 \eqref{BorelIntegral} 的 Borel 积分（Borel 不可求和的 renormalon），必须与非微扰修正相匹配。注意，\eqref{B[H]} 中的 $\widetilde{R}(w)$ 项是 $w$ 的解析函数，因此与奇点的讨论无关。

%
\subsubsection{$w= 1/2 $ 处的紫外 renormalon}\label{sect_uvren}%
%
$w=1/2$ 处的奇点由 Wilson 线自能插入中的大动量贡献生成，它是重正化因子的一部分：
\begin{align}
  B[H] &\stackrel{w\to 1/2}{=} \frac{-4C_F}{w-1/2}\sqrt{X}.
\end{align}
这一奇点众所周知 \cite{Beneke:1994sw}，与 Wilson 线自能中的线性 UV 发散一一对应（另见 (\ref{n-qITD}) 之前的讨论）。采用归一化的 qPDF（\eqref{n-qITD} 或 \eqref{hat-qITD}）即可将其消除，本文不再进一步讨论。

%
\subsubsection{红外 renormalon}\label{sect_irren}%
%
领头的红外 renormalon 奇点位于 $w=1$。我们得到
\begin{align}
\label{w=1}
B[H^\parallel](w) &\stackrel{w\to 1}{=} \frac{-4C_F}{1-w}\Big[\alpha + \bar\alpha \ln \bar\alpha\Big]X\,,
\notag\\
B[H^\perp](w) &\stackrel{w\to 1}{=} \frac{-4C_F}{1-w}\Big[\alpha + \bar\alpha \ln \bar\alpha + \alpha\bar\alpha\Big]X\,.
\end{align}
提醒一下我们的记号 $\bar\alpha = 1-\alpha$。
这两个表达式是我们的主要结果。

$w=n$（$n=2,3...$）处的 renormalon 奇点具有如下一般形式
\begin{align}
B[H](w)&=\frac{2C_F}{n-w}\Big[\alpha^np_{n-1}(\alpha)+\frac{(-1)^n\delta(\bar \alpha)}{n!(n\!-\!2)!n^2(2n\!-\!1)}\Big] X^w,
\end{align}
其中 $p_n(\alpha)$ 是 $n$ 次多项式，例如 $p_1(\alpha)=(5\alpha-3)/6$、$p_2(\alpha)=(\alpha^2-25\alpha+20)/180$ 等。


%
